\section{Causal Judgment: What the Evidence Permits}

The temptation is symmetrical. One side sees a council, a new Mass, and descending graphs, then declares the case closed. The other sees secularization across the West and declares Catholic choices causally irrelevant. Both confuse the presence of multiple causes with the absence of responsibility.

Historical causation is normally conjunctural. External pressure acts through institutions; institutional decisions alter resistance or susceptibility; demographics change the pool; symbols alter identity; scandals destroy accumulated trust. The question is not “Council or culture?” but which mechanisms joined them, with what timing and regional variation.

\subsection{The evidentiary ledger}

\begin{longtable}{>{\bfseries\raggedright\arraybackslash}p{0.27\linewidth}>{\raggedright\arraybackslash}p{0.31\linewidth}>{\raggedright\arraybackslash}p{0.34\linewidth}}
\toprule
\textbf{Proposition} & \textbf{Evidence supporting it} & \textbf{Reason not to overstate it}\\
\midrule
\endhead
The crisis is temporally postconciliar. & Major United States indicators turn sharply after their mid-1960s peak; attendance, vows, ordinations, baptisms, and marriages fall during implementation. & Some trends and cultural causes predate the council; full books arrived after decline had begun.\\
Vatican II authorized destabilizing breadth. & Liturgy, institutes, formation, ecumenism, governance, and pastoral stance were all officially revised. & The texts also command continuity, founder charism, Latin, chant, discipline, doctrine, and holiness; not every implementation follows them.\\
The liturgical reform plausibly mediated identity loss. & Weekly worship is high-frequency formation; nearly every sensory signal changed; replacement was rapid; attendance and belief later weakened. & Same Missal produced different local outcomes; other churches declined; no controlled counterfactual isolates the rite.\\
Religious-life reforms contributed to institute decline. & Special chapters revised core practices during mass departures; identity clarity predicts attraction in some later studies. & The mid-century peak, aging, fertility, and women's opportunities independently reduce numbers; institute outcomes vary.\\
The 1960s cultural revolution was a major cause. & Parallel changes in sex, family, authority, education, and affiliation crossed denominations and countries. & Catholic institutions chose how to teach, discipline, worship, and respond; external pressure does not erase agency.\\
Public dissent weakened the teaching office. & \emph{Humanae vitae} rejection normalized selective assent; later survey attitudes show broad divergence. & Dissent may also be an effect of earlier weak formation and social change; one controversy does not explain every doctrine.\\
Abuse and concealment deepened collapse. & Victimization and institutional failure devastated moral credibility and trust, especially after public revelations. & Much vocation and attendance decline preceded broad disclosure; it is a later accelerator, not the single origin.\\
Global growth refutes a universal collapse. & Catholic population, diocesan priests in count, parishes, deacons, and schools grew worldwide in selected periods. & Growth often lagged population; religious density and sacraments per Catholic deteriorated; regions differ.\\
Traditional communities show an alternative. & Some communities with stable liturgy, doctrine, common life, and clear identity attract young families and vocations. & Small bases, migration, self-selection, fertility, and concentrated commitment prevent simple extrapolation to the whole Church.\\
\bottomrule
\end{longtable}

\subsection{Mechanism 1: authorized boundary uncertainty}

Institutions teach not only propositions but which propositions can change. The council and postconciliar authorities lawfully changed many disciplines and practices while reaffirming doctrine. When leaders failed to explain the boundary, faithful inferred that all inherited forms were provisional. The expectation of further change made present teaching less costly to ignore: a person could dissent while waiting for official revision.

This mechanism does not require a false conciliar sentence. It requires simultaneous visible change, contested expert commentary, weak catechesis, and delayed correction. \emph{Humanae vitae} made the boundary crisis explicit. The Church said one moral norm remained; influential Catholics said the process of development had outrun the judgment. Selective assent became socially organized.

\subsection{Mechanism 2: loss of redundant transmission}

Preconciliar Catholicism transmitted identity through overlapping systems: Sunday Mass, Friday abstinence, Ember days, fasts, sodalities, processions, schools staffed by religious, devotions, parish boundaries, clerical dress, family size, and neighborhood institutions. Not every element was essential; together they made Catholic difference difficult to miss.

Many changed because of urbanization, affluence, ethnicity, mobility, and professionalization. Ecclesial reform removed or weakened others at the same time. The loss was not merely a list of rules. Redundancy protects transmission: when one channel fails, another carries the pattern. Simultaneous simplification can make doctrine depend on explicit catechesis precisely when schools and families are weakening.

\subsection{Mechanism 3: ritual discontinuity and memory}

Liturgy joins generations through repeated words and gestures. Rapid replacement can tell older faithful that their religious memory is obsolete and younger faithful that worship is a designed product. If celebrants visibly choose and improvise, the rite may appear to originate in the community rather than be received. That perception can weaken transcendent authority even if the General Instruction teaches the contrary.

The mechanism is plausible, not universal. Vernacular proclamation can deepen reception; careful celebration can manifest continuity; some older practices obscured meaning. The causal claim is conditional: where reform was rapid, iconoclastic, variable, and poorly catechized, it likely accelerated loss of sacred and institutional continuity.

\subsection{Mechanism 4: altered incentives and institute ecology}

Religious life had offered a thick common rule, clear status, sacrificial purpose, and institutional mission. Wider professional opportunity reduced practical incentives. When institutes simultaneously relaxed common structures and recast authority, the vocational cost remained—poverty, chastity, obedience—while distinctive communal rewards could appear lower. Rodney Stark and Roger Finke formulate this as a cost-and-reward model; Patricia Wittberg analyzes religious orders as social movements sustained by ideology, ritual, recruitment, and organization. Their theories do not exhaust supernatural vocation, but they identify observable institutional conditions.

Stephen Bullivant's account of Catholic disaffiliation emphasizes a “perfect storm”: pre-existing social change, the council's upheaval and expectations, the sexual revolution, and weak transmission interacted. That conjunctural model fits the evidence better than either a self-contained conciliar cause or an external secularization that leaves Catholic leaders causally innocent.

\subsection{Tests that complicate monocausality}

\textbf{Other churches declined.} Mainline Protestant communities in similar cultures lost attendance and affiliation without Vatican II or the Roman Missal. This strongly supports common secularizing causes. It does not show Catholic reform had zero marginal effect.

\textbf{One reform, different outcomes.} The same Latin typical Missal underlies growing African parishes, collapsing European ones, reverent monasteries, and improvisational suburban celebrations. Demography, clergy, culture, translation, and local implementation mediate the result. The book cannot be a sufficient cause.

\textbf{One culture, different Catholic subgroups.} Within Western societies, communities marked by frequent worship, doctrinal clarity, common prayer, dense family networks, or traditional liturgy sometimes retain more strongly. Selection matters: already committed people choose demanding communities. But selection itself shows that costly clarity can sustain identity rather than necessarily repel.

\textbf{Timing differs by measure.} Women religious peak and decline differently from infant baptisms, school enrollment, priest counts, abuse disclosure, and self-identification. A single switch cannot explain staggered mechanisms.

\subsection{The unavailable counterfactuals}

No historian observes the same 1960s Catholic population both with and without Vatican II. Four counterfactuals are nevertheless useful if held conditionally.

\begin{itemize}
\item \textbf{No council:} Western secularization, fertility decline, sexual revolution, and mobility would still have pressured Catholic institutions. Some decline was highly likely. An unreformed defensive system might have retained stronger boundaries for a time while failing differently in mission and credibility.
\item \textbf{Council without rapid comprehensive implementation:} slower liturgical change, parallel availability of inherited books, stronger doctrinal boundary-setting, and institute-specific renewal might plausibly have reduced disruption. This is not provable, but the mechanism is coherent.
\item \textbf{Missal reform with disciplined continuity:} Latin and chant retained in practice, fewer options, stable orientation and architecture, faithful translation, and rigorous catechesis might have realized conciliar goods with less identity loss. Some communities approximate this, but selection prevents a clean experiment.
\item \textbf{No conciliar reform but the same abuse and cultural failures:} trust and practice would still have suffered severely. Liturgy cannot compensate for unrepented institutional injustice or absent family transmission.
\end{itemize}

Counterfactual modesty cuts both ways. Defenders cannot assert that decline would certainly have been worse without the council. Critics cannot assert that 1950s practice would have continued indefinitely with the 1962 Missal.

\subsection{A proportionate verdict}

The council was not a sociological vaccine and should not be defended by claims that its “real fruits” will appear after every contrary generation passes. Neither was it a program to empty convents and deny the Eucharist. Its texts contain resources for renewal: holiness, Scripture, liturgy, founder charism, mission, episcopal responsibility, and serious engagement with atheism. They also authorized a breadth of change whose postconciliar management often underestimated how institutions carry faith.

Paul VI's Missal was not a doctrinal poison. Its reform was a high-salience institutional intervention that changed the ordinary vehicle of Catholic memory. Given the timing, scale, and mechanisms, it is reasonable to judge that the reform as distributed and locally implemented contributed to and accelerated Western disaffiliation in some measure. Available data cannot quantify that measure or separate the official book from surrounding change.

Catholic leadership bears real responsibility where it confused doctrine and adaptation, tolerated dissent, trivialized sacred signs, dismantled effective practices without replacements, or hid grave sin. That responsibility operates beneath the level of conciliar validity. One need not deny the Holy Spirit's protection of the Church to admit that protected authorities can govern imprudently and subordinates can abuse lawful mandates.

\judgmentbox{The best-supported causal model is a convergence: long Western secularization, demographic and economic change, the sexual revolution, council-authorized institutional revision, rapid liturgical reconstruction and replacement, weak boundary catechesis, organized dissent, institute-level choices, and later abuse revelations interacted. Vatican II and Paul VI's Missal were neither necessary nor sufficient causes of the whole crisis. Their authorized implementation was a material mediator and likely accelerator in important Western settings.}
